\section{Numerical and inference validation}
\label{app:numerical-validation}

The test suite includes the Einstein--de Sitter expansion law, a sparse positive-redshift distance regression, monotonic luminosity distance, BAO observable definitions, positive-definite DESI covariance validation, an exact flat de Sitter scalar-field limit, potential-derivative verification, analytical bounce and turnaround cases, topology-aware return metrics, a constant-field $S^4$ Laplace--Beltrami null mode, finite-volume layered-field energy conservation, and manifest-schema validation. The sparse-redshift test specifically prevents the first DESI distance from being assigned zero merely because the requested measurement grid begins above $z=0$.

The background verification workflow executes the declared tolerance ladder and independent-integrator comparison described in the main text. Solver verification always compares the unwrapped scalar trajectory. A displacement by one or more periods of the potential is retained because a periodic potential alone does not identify the scalar target space. Turning states come from root-localized adaptive output and are serialized separately from the uniformly sampled plotting trajectory, so return metrics do not depend on the requested plotting density. Even repeated close returns in one integration are not a proof of recurrence or a stable limit cycle; Poincar\'e maps, variational equations, Floquet multipliers, basin studies, perturbation stability, and synthetic recovery remain later-stage gates.

\subsection{DESI chain convergence and posterior construction}
The DESI inference workflow downloads the official sample files and release-side metadata, including the checkpoint and updated configuration, and records SHA-256 checksums. GetDist independently evaluates the same sample files used for the plotted posterior. The burn-selection rule is predeclared: test candidate initial-row fractions and select the earliest cut satisfying multivariate $R-1\le0.01$. A command-line override is accepted only if the same convergence threshold still passes.

\IfFileExists{generated/desi_chain_convergence_table.tex}{
\begin{table}[ht]
\centering
\small
\input{generated/desi_chain_convergence_table.tex}
\caption{GetDist convergence diagnostics for the official DESI chain sets. Weighted ESS here is the standard weight-only effective sample count $(\sum_iw_i)^2/\sum_iw_i^2$; the serialized GetDist convergence report additionally includes autocorrelation and split-chain diagnostics.}
\label{tab:desi-convergence}
\end{table}}{}

The principal two-dimensional regions are produced from GetDist's marginalized density estimator rather than a hand-tuned image-space Gaussian smoothing. This makes the contour construction auditable against a standard cosmological-chain analysis package \cite{getdist}. The posterior-predictive panel uses the complete pinned covariance through its Cholesky factor $C=LL^T$ and plots $L^{-1}(m-d)$ rather than independent marginal residuals.

\subsection{Derived kinematics and sound-horizon cross-check}
The late-time flat-$\Lambda$CDM acceleration quantities are propagated through the released BAO+BBN posterior rather than shown as exact central-model curves.
\IfFileExists{generated/kinematic_summary_table.tex}{
\begin{table}[ht]
\centering
\small
\input{generated/kinematic_summary_table.tex}
\caption{Posterior-derived flat-$\Lambda$CDM kinematic summaries. These quantities are standard-cosmology consequences of the released posterior, not independent observables.}
\label{tab:kinematics}
\end{table}}{}

The Aubourg approximation is not the inference backend. It is evaluated over a deterministic posterior-spanning set and compared directly with CAMB.
\IfFileExists{generated/aubourg_validation_table.tex}{
\begin{table}[ht]
\centering
\small
\input{generated/aubourg_validation_table.tex}
\caption{Sound-horizon backend cross-check over posterior-spanning parameter draws.}
\label{tab:aubourg-validation}
\end{table}}{}
\IfFileExists{generated/sound_horizon_backend_validation.png}{
\begin{figure}[ht]
\centering
\includegraphics[width=0.72\linewidth]{generated/sound_horizon_backend_validation.png}
\caption{Fractional difference between the Aubourg sound-horizon approximation and CAMB. The approximation remains a supplementary numerical diagnostic only.}
\label{fig:aubourg-validation}
\end{figure}}{}
